% 見開きのダイジェスト。表紙の次の 2 ページ（左が介入、右が検査）に置く。
\newcommand{\dghead}[2]{%
  \noindent{\kkgt{4}\small\color{kkaccent}#1}\par\vspace{1mm}%
  \noindent{\kkgt{7}\fontsize{17pt}{21pt}\selectfont #2}\par\vspace{1.5mm}%
  \noindent{\color{kkaccent}\rule{\linewidth}{1pt}}\par\vspace{2mm}}
\newcommand{\dgstep}[3]{%
  \par\noindent\begin{minipage}[t]{6.5mm}\vspace{0pt}\centering
    \tikz\node[circle, fill=kkaccent, text=white, inner sep=0pt,
      minimum size=4.6mm, font=\scriptsize\bfseries]{#1};\end{minipage}%
  \begin{minipage}[t]{\dimexpr\linewidth-6.5mm\relax}\vspace{0pt}
    \small\linespread{1}\selectfont{\kkgt{5}\color{kkdark}#2}\quad #3\end{minipage}\par\vspace{2mm}}

\newgeometry{top=12mm,bottom=12mm,left=14mm,right=14mm}
\thispagestyle{empty}
\dghead{ダイジェスト 1}{介入のメタ解析の手順}

{\small メタ解析の手順は、治療や予防などの介入の効果をまとめる方法として整えられてきました。検査のメタ解析も工程はほぼ同じで、主に統計の部分が変わります。そこで、基本になる介入の手順をこのページに、検査で変わるところを次のページに並べました。初めは用語が分からなくても、眺めるだけで構いません。各段階の末尾の括弧は、詳しく書いた章です。}
\vspace{2mm}

\dgstep{1}{問いを PICO で書く}{対象、介入、比較、アウトカムを一文に書きます。主要アウトカムを一つに決めます（第 2 章）。}
\dgstep{2}{解析より先に計画を書いて登録する}{効果指標、モデル、サブグループ、感度分析を決めてから PROSPERO に登録します。結果を見て解析を選ぶと結論が偏るからです（第 2 章）。}
\dgstep{3}{複数のデータベースと試験登録を検索する}{同義語を OR、概念どうしを AND でつなぎます。未発表の試験は試験登録で探します（第 3 章）。}
\dgstep{4}{2 人で独立に選別する}{抄録で除外し、残りは本文で判断します。除外の理由と件数を PRISMA のフロー図に残します（第 3 章）。}
\dgstep{5}{数を抜き出し、バイアスを評価する}{群ごとの事象数と人数を 2 人で抜き出し、RoB 2 で試験ごとに判定します（第 4 章）。}
\dgstep{6}{効果を統合する}{リスク比を対数にして、精度の高い試験ほど重く平均します。既定はランダム効果モデル（REML）にして、試験数が少なければ、区間を広めに作る HKSJ 法の区間も並べて比べます（第 5 章）。}
\dgstep{7}{ばらつきを読む}{I² だけで判断せず、τ² と予測区間を見ます。サブグループの差は試験間の比較にすぎないと読みます（第 6 章）。}
\dgstep{8}{結果を疑う}{ファネルプロット、バイアスのリスクの高い試験を除いた感度分析を行い、GRADE で確実性を判定します（第 7 章）。}
\dgstep{9}{絶対効果に直して報告する}{1000 人あたり何人の差になるかを示し、PRISMA 2020 に沿って書きます（第 7、8 章）。}

\vspace{1mm}
{\small R での計算は \kkref{sec:05-6} で示します。架空の 10 試験では、リスク比 0.73（95\% 信頼区間 0.61〜0.86）、I² 51\%、予測区間 0.46〜1.15 でした。}

\clearpage
\thispagestyle{empty}
\dghead{ダイジェスト 2}{検査のメタ解析で変わるところ}

{\small 工程は前のページとほぼ同じです。大きく違うのは、結果が感度と特異度の 2 つの数で、互いに関係していることです。}
\vspace{2mm}

\dgstep{1}{問いを 4 つの要素で書く}{対象、対象検査、参照基準、標的疾患を決め、検査を臨床のどこで使うのか（置き換え、トリアージ、追加）を書きます（第 11 章）。}
\dgstep{2}{検索にフィルターをかけない}{診断精度研究を確実に拾う検索フィルターはないので、検査名と疾患名だけで探します（第 11 章）。}
\dgstep{3}{研究ごとに 2×2 表を作る}{真陽性、偽陽性、偽陰性、真陰性の 4 つの数を抜き出し、カットオフ、測定試薬、検体の違いも記録します（第 12 章）。}
\dgstep{4}{QUADAS で評価する}{患者の選び方、対象検査、参照基準、フローとタイミングの 4 領域で評価します。症例対照デザインや一部の人だけの検証は、精度を高く見せます。2026 年に改訂版の QUADAS-3 が出ています（第 12 章）。}
\dgstep{5}{カットオフをそろえて統合する}{要約の感度と特異度は、同じカットオフの研究だけで出します。2 つを同時に扱う二変量モデルを使い、カットオフの混ざった研究全体には SROC 曲線を描きます（第 10、13 章）。}
\dgstep{6}{ばらつきの原因を探す}{検査のメタ解析ではばらつきが大きいのが普通です。カットオフ、場面、測定法を共変量にして調べ、予測領域で次の研究がどこに来るかを見ます（第 14 章）。}
\dgstep{7}{臨床の言葉にして報告する}{有病割合ごとに 1000 人あたりの真陽性、偽陽性、偽陰性、真陰性を示し、PRISMA-DTA に沿って書きます（第 15 章）。}

\vspace{1mm}
{\small R での計算は \kkref{sec:13-3} で示します。細菌感染症を見分けるプロカルシトニン（PCT）の架空の 10 研究のうち、カットオフ 0.5 ng/mL の 6 研究では、感度 0.75、特異度 0.76、陽性尤度比 3.1、陰性尤度比 0.34 でした。有病割合 10\% なら、陽性と出た人のうち菌血症は 4 人に 1 人程度です。}

\vspace{1mm}
\noindent{\footnotesize\color{kkdark} 指針の使い分け: 個々の診断精度研究の報告は STARD 2015、検査のメタ解析の報告は PRISMA-DTA、研究ごとのバイアスの評価は QUADAS-2（または QUADAS-3）を使います。}
\clearpage
\restoregeometry
